\section{Kantorovich Relaxation}
\label{sec-kantorovich}
\subsection{Discrete Relaxation}
\label{sec-discrete-relaxation}
\paragraph{Mass splitting and couplings.}
\(\al=\sum_i \a_i\de_{x_i}, \qquad \be=\sum_j \b_j\de_{y_j}.\)
\begin{defn}[Discrete couplings and mass conservation]\label{def-discrete-couplings}
Admissible couplings are only constrained to satisfy conservation of mass:
\eql{\label{eq-discr-couplings}
\CouplingsD(\a,\b)
\eqdef
\enscond{\P\in\RR_+^{n\times m}}{\P\ones_m=\a \qandq \transp{\P}\ones_n=\b}.
}
Equivalently, \(\P\ones_m=\bigl(\sum_j \P_{i,j}\bigr)_i\in\RR^n, \qquad \transp{\P}\ones_n=\bigl(\sum_i \P_{i,j}\bigr)_j\in\RR^m.\)
\end{defn}
\begin{defn}[Discrete product coupling]\label{def-discrete-product-coupling}
Given weights $\a\in\simplex_n$ and $\b\in\simplex_m$, the discrete product, or trivial, coupling is \((\a\otimes\b)_{i,j}\eqdef \a_i\b_j.\)
It belongs to $\CouplingsD(\a,\b)$ and corresponds to choosing the source and target labels independently.
\end{defn}
\begin{prop}[Discrete product optimality is degenerate]\label{prop-discrete-product-coupling-degenerate}
Assume that the zero-mass rows and columns have been removed, so that $\a_i>0$ and $\b_j>0$, and let $\C$ be a finite cost matrix. The product plan $\a\otimes\b$ minimizes $\P\mapsto\dotp{\C}{\P}$ over $\CouplingsD(\a,\b)$ if and only if every coupling $\P\in\CouplingsD(\a,\b)$ minimizes it.
\end{prop}
\begin{proof}
The reverse implication is immediate. Assume conversely that $\a\otimes\b$ is optimal and let $\Q\in\CouplingsD(\a,\b)$ be arbitrary. Since all entries of $\a\otimes\b$ are positive, there exists $t>0$ small enough that
\(\R\eqdef (1+t)(\a\otimes\b)-t\Q\)
has nonnegative entries. Its row and column sums are $\R\ones_m=(1+t)\a-t\a=\a$ and $\R^\top\ones_n=(1+t)\b-t\b=\b$, so $\R\in\CouplingsD(\a,\b)$. Moreover,
\(\a\otimes\b=\frac{1}{1+t}\R+\frac{t}{1+t}\Q\).
By optimality of $\a\otimes\b$, both $\dotp{\C}{\R}$ and $\dotp{\C}{\Q}$ are at least $\dotp{\C}{\a\otimes\b}$. Taking the scalar product of the convex-combination identity with $\C$ gives
\(\dotp{\C}{\a\otimes\b} = \frac{1}{1+t}\dotp{\C}{\R} +\frac{t}{1+t}\dotp{\C}{\Q}.\)
A convex average of two numbers not smaller than $\dotp{\C}{\a\otimes\b}$ can equal $\dotp{\C}{\a\otimes\b}$ only if both numbers are equal to it. Hence $\dotp{\C}{\Q}=\dotp{\C}{\a\otimes\b}$, and $\Q$ is optimal. Since $\Q$ was arbitrary, all couplings are optimal.
\end{proof}
Thus the product plan is mainly a feasibility witness. Except in the degenerate situation where the linear cost is constant on the whole transportation polytope, it is not expected to solve optimal transport.

\paragraph{Linear-programming structure.}
\begin{defn}[Discrete Kantorovich problem]\label{def-discrete-kantorovich-problem}
For $\a\in\simplex_n$, $\b\in\simplex_m$, and $\C\in\RR^{n\times m}$, define
\eql{\label{eq-kanto-discr}
\MKD_{\C}(\a,\b)
\eqdef
\umin{\P\in\CouplingsD(\a,\b)}
\dotp{\C}{\P}
=
\umin{\P\in\CouplingsD(\a,\b)}
\sum_{i,j}\C_{i,j}\P_{i,j}.
}
A minimizer $\P$ is an optimal coupling, or optimal transport plan.
\end{defn}
\begin{prop}[Convexity in the marginals and concavity in the cost]\label{prop-discrete-kantorovich-joint-convexity}
For every fixed cost matrix $\C\in\RR^{n\times m}$, the map
\((\a,\b)\longmapsto\MKD_{\C}(\a,\b)\)
is jointly convex on $\simplex_n\times\simplex_m$. For every fixed $(\a,\b)\in\simplex_n\times\simplex_m$, the map
\(\C\longmapsto\MKD_{\C}(\a,\b)\)
is concave on $\RR^{n\times m}$.
\end{prop}
\begin{proof}
For $r\in\{0,1\}$, let $\P_r$ be optimal between $\a_r$ and $\b_r$. For $t\in[0,1]$, the matrix $\P_t=(1-t)\P_0+t\P_1$ belongs to $\CouplingsD((1-t)\a_0+t\a_1,(1-t)\b_0+t\b_1)$. Therefore
\(\MKD_{\C}\big((1-t)\a_0+t\a_1,(1-t)\b_0+t\b_1\big) \leq \dotp{\C}{\P_t} = (1-t)\MKD_{\C}(\a_0,\b_0)+t\MKD_{\C}(\a_1,\b_1).\)
For fixed marginals, every $\P\in\CouplingsD(\a,\b)$ satisfies \(\dotp{(1-t)\C_0+t\C_1}{\P} \geq (1-t)\MKD_{\C_0}(\a,\b) +t\MKD_{\C_1}(\a,\b).\)
Minimizing the left-hand side over $\P$ proves concavity in $\C$.
\end{proof}
\begin{prop}[Rank-controlled sparse minimizers]\label{prop-lp-rank-sparsity}
Let $\mathcal A:\RR^{n\times m}\to\RR^q$ be linear, let $\mathbf d\in\RR^q$ and $\C\in\RR^{n\times m}$, and consider
\begin{equation}\label{eq-lp-rank-sparsity}
\inf_{\P\geq 0,\;\mathcal A(\P)\leq\mathbf d}
\dotp{\C}{\P},
\end{equation}
where the inequality is componentwise. If the feasible set is nonempty and the objective is bounded below on it, then~\eqref{eq-lp-rank-sparsity} admits a minimizer $\P^\star$ satisfying
\begin{equation}\label{eq-lp-rank-support-bound}
\#\enscond{(i,j)}{\P^\star_{i,j}>0}
\leq \operatorname{rank}(\mathcal A).
\end{equation}
Here $\operatorname{rank}(\mathcal A)=\dim(\operatorname{Im}\mathcal A)$. The same conclusion holds when $\mathcal A(\P)\leq\mathbf d$ is replaced by $\mathcal A(\P)=\mathbf d$.
\end{prop}
\begin{proof}
The feasible region is a nonempty polyhedron, and the objective is bounded below on it. The attainment theorem for linear programs therefore provides a minimizer. Among all minimizers, choose $\P^\star$ with the smallest number of positive entries, and denote its support by $S$.

Suppose that $\#S>\operatorname{rank}(\mathcal A)$. The restriction of $\mathcal A$ to the $\#S$-dimensional space of matrices supported on $S$ then has a nontrivial kernel. Hence there is a nonzero matrix $H$, supported on $S$, such that $\mathcal A(H)=0$. Since $\P^\star$ is strictly positive on $S$, both $\P^\star+tH$ and $\P^\star-tH$ remain nonnegative for all sufficiently small $t>0$. Moreover, $\mathcal A(\P^\star\pm tH)=\mathcal A(\P^\star)$, so both perturbations satisfy exactly the same linear inequalities, or equalities, as $\P^\star$. Optimality forces $\dotp{\C}{H}=0$; otherwise one of the two signs would strictly decrease the objective.

Choose $\sigma\in\{-1,1\}$ such that $\sigma H$ has a negative entry, and set $t_\star\eqdef\min_{(i,j):\,\sigma H_{i,j}<0}\P^\star_{i,j}/(-\sigma H_{i,j})>0$.
Then $\P^\star+t_\star\sigma H$ is nonnegative, feasible, and has the same cost as $\P^\star$. At least one of its positive entries has become zero, while no new entry has appeared outside $S$. This contradicts the minimality of $S$ and proves~\eqref{eq-lp-rank-support-bound}.
\end{proof}
\begin{prop}[Sparse optimal plans]\label{prop-sparse-optimal-plans}
Assume $\a_i>0$, $\b_j>0$ and $\sum_i \a_i=\sum_j \b_j=1$. The linear program~\eqref{eq-kanto-discr} admits an optimal coupling with at most $n+m-1$ nonzero entries.
\end{prop}
\begin{proof}
The transportation polytope is nonempty and compact, so the minimum is finite. Apply the equality case of Proposition~\ref{prop-lp-rank-sparsity} to the marginal operator $\mathcal A_{\mathrm{marg}}(\P)\eqdef(\P\ones_m,\transp{\P}\ones_n)\in\RR^{n+m}$.
Its rank is $n+m-1$. Indeed, every matrix satisfies
$\sum_i(\P\ones_m)_i=\sum_j(\transp{\P}\ones_n)_j$, so there is one linear dependence among the $n+m$ marginal coordinates. It is the only one: if $u\in\RR^n$ and $v\in\RR^m$ satisfy
\(\dotp{u}{\P\ones_m} + \dotp{v}{\transp{\P}\ones_n}=0 \qquad\text{for every }\P,\)
then testing this identity on each elementary matrix gives $u_i+v_j=0$ for every $(i,j)$. Thus all $u_i$ equal one scalar and all $v_j$ equal its opposite, so the annihilator of $\operatorname{Im}(\mathcal A_{\mathrm{marg}})$ is one-dimensional. It follows that $\operatorname{rank}(\mathcal A_{\mathrm{marg}})=n+m-1$, and Proposition~\ref{prop-lp-rank-sparsity} supplies the claimed optimal coupling.
\end{proof}
\begin{prop}[North-west corner feasible plan]\label{prop-northwest-corner}
Let $\a\in\RR_+^n$ and $\b\in\RR_+^m$ have the same positive total mass. Starting from $(i,j)=(1,1)$ with residual masses $r_i=\a_i$ and $s_j=\b_j$, skip zero residuals, set $\P_{i,j}=\min(r_i,s_j)$,
subtract this value from both residuals, and advance every index whose residual has become zero. Repeat until all residual masses are exhausted.
\end{prop}
\begin{proof}
All assignments are nonnegative. At each step, the mass placed in entry $(i,j)$ is subtracted from exactly one current row residual and one current column residual, so no row or column can receive more mass than prescribed. Conversely, an index is advanced only when its residual has been fully filled. When the algorithm stops, the total assigned mass is $\sum_i \a_i=\sum_j \b_j$, hence all row and column sums are exactly $\a$ and $\b$.

Each positive assignment exhausts at least one current row or one current column. Before the final assignment, at most $n-1$ row advances and $m-1$ column advances can occur without terminating the construction. Hence the number of positive entries is at most $(n-1)+(m-1)+1=n+m-1$. For acyclicity, view the positive support as a bipartite graph. Once a row or column index is advanced, it never appears again, so each new positive edge either starts a new component or attaches at least one new vertex to the component currently being swept. No edge is ever added between two old vertices of the same component, so no cycle can be created.
\end{proof}
\paragraph{One-dimensional cases.}
\begin{prop}[One-dimensional weighted sweep]\label{prop-1d-weighted-sweep}
Let $x_1\leq\cdots\leq x_n$ and $y_1\leq\cdots\leq y_m$ be points on the line, and let $c(x,y)=h(x-y)$ with $h$ convex. The north-west corner plan between the sorted weighted atoms is optimal for~\eqref{eq-kanto-discr}.
\end{prop}
\begin{proof}
The north-west plan is monotone: if $i<i'$ and $j>j'$, it cannot put positive mass on both $(i,j)$ and $(i',j')$, because the sweep exhausts rows and columns in increasing order. To prove optimality, start from an arbitrary feasible plan $\P$. If
$\P_{1,1}<\min(\a_1,\b_1)$, then row $1$ sends positive mass to some column $j>1$, while column $1$ receives positive mass from some row $i>1$. Moving
\(\eta=\min(\P_{1,j},\P_{i,1})\)
from $(1,j)$ and $(i,1)$ to $(1,1)$ and $(i,j)$ preserves both marginals. Convexity of $h$ gives the Monge inequality
\(h(x_1-y_j)+h(x_i-y_1) \geq h(x_1-y_1)+h(x_i-y_j),\)
so this exchange does not increase the cost. Each exchange either removes the chosen entry in row $1$ or removes the chosen entry in column $1$. Repeating finitely many times forces
$\P_{1,1}=\min(\a_1,\b_1)$, exactly the first north-west assignment. It exhausts row $1$, column $1$, or both. Removing every exhausted index reduces the problem to a smaller ordered transportation problem. Induction on $n+m$ produces the complete north-west plan without increasing cost. Sorting costs $O(n\log n+m\log m)$ and the sweep uses at most $n+m-1$ positive assignments.
\end{proof}
\paragraph{Permutation matrices as couplings.}
\begin{defn}[Permutation matrices]\label{def-permutation-matrices}
For a permutation $\sigma\in\Perm(n)$, its permutation matrix $\P_\sigma$ is
\[
(\P_\sigma)_{i,j}=\choice{1 \qifq j=\sigma(i)\\ 0 \quad\text{otherwise}.}
\]
The set of all permutation matrices is \(\mathcal P_n^{\mathrm{perm}} \eqdef \enscond{\P_\sigma}{\sigma\in\Perm(n)}.\)
\end{defn}
\(\dotp{\C}{\P_\sigma/n}=\frac1n\sum_{i=1}^n \C_{i,\sigma(i)}.\)
\begin{defn}[Birkhoff polytope]\label{def-birkhoff-polytope}
The Birkhoff polytope is the convex set of bistochastic matrices
\(\mathcal B_n \eqdef \enscond{\P\in\RR_+^{n\times n}}{\P\ones_n=\ones_n \qandq \transp{\P}\ones_n=\ones_n}.\)
\end{defn}
Then $\CouplingsD(\ones_n/n,\ones_n/n)=\mathcal B_n/n$, and permutation couplings are included in this convex relaxation. More precisely,

\(\mathcal P_n^{\mathrm{perm}} = \mathcal B_n\cap\{0,1\}^{n\times n} \subset \mathcal B_n,\)
\(\min_{\P\in\mathcal B_n}\dotp{\C}{\P} \leq \min_{\P\in\mathcal P_n^{\mathrm{perm}}}\dotp{\C}{\P}.\)
\begin{defn}[Extreme points]\label{def-extreme-points}
For a convex set $\mathcal C$ in a finite-dimensional vector space, \(\operatorname{Extr}(\mathcal C) \eqdef \enscond{x\in\mathcal C}{x=(y+z)/2,\ y,z\in\mathcal C \Rightarrow y=z=x}.\)
\end{defn}
\begin{prop}[Existence of extreme points]\label{prop-extreme-point-existence}
If $\mathcal C$ is a non-empty compact convex subset of a finite-dimensional vector space, then $\operatorname{Extr}(\mathcal C)$ is non-empty.
\end{prop}
\begin{proof}
Among all non-empty faces of $\mathcal C$, choose one of minimal affine dimension. If this face contained two distinct points, maximizing a linear functional that is not constant on the face would produce a non-empty proper exposed subface, contradicting minimality. Hence the minimal face is a singleton, and its point is extreme.
\end{proof}
\begin{prop}[Linear programs have extreme minimizers]\label{prop-linear-program-extreme-minimizer}
Let $\mathcal C$ be a non-empty compact convex set. For every linear form $\ell$, \(\operatorname{Extr}(\mathcal C)\cap\operatorname*{argmin}_{x\in\mathcal C}\ell(x) \neq\emptyset.\)
\end{prop}
\begin{proof}
The set $S=\operatorname*{argmin}_{x\in\mathcal C}\ell(x)$ is non-empty, compact and convex. By Proposition~\ref{prop-extreme-point-existence}, it has an extreme point $x$. If $x=(y+z)/2$ with $y,z\in\mathcal C$, then by linearity and optimality of $x$, both $y$ and $z$ also minimize $\ell$ on $\mathcal C$, hence $y,z\in S$. Since $x$ is extreme in $S$, $y=z=x$. Thus $x$ is extreme in $\mathcal C$.
\end{proof}
\begin{thm}[Birkhoff--von Neumann]\label{thm-birkhoff-von-neumann}
The extreme points of $\mathcal B_n$ are exactly the permutation matrices.
\end{thm}
\begin{proof}
We first prove that permutation matrices are extreme. Let $\P_\sigma\in\mathcal P_n^{\mathrm{perm}}$ and assume that $\P_\sigma=(\Q+\R)/2$ with $\Q,\R\in\mathcal B_n$.
Every bistochastic matrix has entries in $[0,1]$. Since the only extreme points of $[0,1]$ are $0$ and $1$, each entry of $\P_\sigma$ fixes the corresponding entries of $\Q$ and $\R$: if $(\P_\sigma)_{i,j}=0$, then $\Q_{i,j}=\R_{i,j}=0$, while if $(\P_\sigma)_{i,j}=1$, then $\Q_{i,j}=\R_{i,j}=1$. Hence $\Q=\R=\P_\sigma$, so $\P_\sigma$ is extreme.

Pick $\P\in\mathcal B_n\setminus\mathcal P_n^{\mathrm{perm}}$. Since an integral bistochastic matrix is necessarily a permutation matrix, $\P$ has at least one fractional entry. We shall split
\(\P=\frac{\Q+\R}{2}\)
with $\Q,\R\in\mathcal B_n$ and $\Q\neq \R$, proving that $\P$ is not extreme.

Associate with $\P$ the bipartite graph whose left vertices are the rows, whose right vertices are the columns, and whose edges are the fractional entries
\(0<\P_{i,j}<1\).
An entry equal to $1$ uses the whole mass of its row and column, so it is isolated in the positive support and does not appear in this fractional graph. If a left vertex $i$ is incident to a fractional edge $(i,j_1)$, then it must be incident to at least one other fractional edge. Indeed, the row sum is one; after the contribution $\P_{i,j_1}\in(0,1)$, a positive amount $1-\P_{i,j_1}$ remains in the same row, and it cannot be carried by an entry equal to $1$. The same argument applies to right vertices, using the column sums. Thus every non-isolated vertex of the fractional graph has degree at least two.

Starting from any fractional edge, one may therefore walk through adjacent fractional edges without immediately backtracking and without getting stuck. Since the graph is finite, some vertex is eventually visited twice; the portion of the walk between the two visits contains a cycle. Choose a shortest such cycle and write it in alternating form $(i_1,j_1,i_2,j_2,\ldots,i_p,j_p)$, with $i_{p+1}=i_1$, where both $(i_s,j_s)$ and $(i_{s+1},j_s)$ are fractional for every $s$. The minimality of the cycle implies that the vertices $i_s$ are all distinct and that the vertices $j_s$ are all distinct. In particular,
\(0<\P_{i_s,j_s}<1 \qandq 0<\P_{i_{s+1},j_s}<1.\)
Define
\(\epsilon \eqdef \min_{1\leq s\leq p} \bigl\{ \P_{i_s,j_s},\, \P_{i_{s+1},j_s},\, 1-\P_{i_s,j_s},\, 1-\P_{i_{s+1},j_s} \bigr\}.\)
All these numbers are positive, so $\epsilon>0$. Split the cycle edges into the two alternating families
\(A\eqdef \{(i_s,j_s)\}_{s=1}^p, \qquad B\eqdef \{(i_{s+1},j_s)\}_{s=1}^p.\)
\[
\Q_{i,j}\eqdef
\begin{cases}
\P_{i,j}, & (i,j)\notin A\cup B,\\
\P_{i,j}+\epsilon/2, & (i,j)\in A,\\
\P_{i,j}-\epsilon/2, & (i,j)\in B,
\end{cases}
\qquad
\R_{i,j}\eqdef
\begin{cases}
\P_{i,j}, & (i,j)\notin A\cup B,\\
\P_{i,j}-\epsilon/2, & (i,j)\in A,\\
\P_{i,j}+\epsilon/2, & (i,j)\in B.
\end{cases}
\]
By the definition of $\epsilon$, all modified entries stay in $[0,1]$, so $\Q$ and $\R$ are nonnegative. Each row vertex $i_s$ of the cycle is incident to exactly one edge of $A$ and one edge of $B$; the $+\epsilon/2$ and $-\epsilon/2$ perturbations therefore cancel in that row. The same cancellation holds in each column vertex $j_s$, and all other rows and columns are unchanged. Finally, $\Q\neq \R$ because $\epsilon>0$ and the cycle is non-empty, while by construction $\P=(\Q+\R)/2$. Thus $\P$ is not extreme.
\end{proof}
\begin{cor}[Kantorovich for matching]\label{cor-kantorovich-matching}
If $m=n$ and $\a=\b=\ones_n/n$, then the discrete Kantorovich problem~\eqref{eq-kanto-discr} admits an optimal solution of the form $\P_\sigma/n$. The associated permutation $\sigma$ solves the assignment problem of Section~\ref{sec-monge-pbm}.
\end{cor}
\begin{proof}
The feasible set is $\mathcal B_n/n$. By Proposition~\ref{prop-linear-program-extreme-minimizer}, the linear objective has an optimal extreme point. Since scaling preserves extreme points and Theorem~\ref{thm-birkhoff-von-neumann} identifies the extreme points of $\mathcal B_n$, this optimizer is $\P_\sigma/n$ for some permutation $\sigma$. Its cost is exactly $n^{-1}\sum_i \C_{i,\sigma(i)}$, so $\sigma$ is an optimal assignment.
\end{proof}
\paragraph{Constructive Birkhoff--von Neumann decomposition.}
\begin{cor}[Birkhoff--von Neumann decomposition]\label{cor-birkhoff-von-neumann-decomposition}
Every $\P\in\mathcal B_n$ admits a representation
\(\P=\sum_{r=1}^{L}\lambda_r\P_{\sigma_r}, \qquad \lambda_r>0, \qquad \sum_{r=1}^{L}\lambda_r=1, \qquad L\leq (n-1)^2+1,\)
where each $\P_{\sigma_r}$ is a permutation matrix.
\end{cor}
\begin{proof}
Let $\mathcal K$ be the convex hull of the permutation matrices. Clearly $\mathcal K\subset\mathcal B_n$. Conversely, if some $\P\in\mathcal B_n$ did not belong to the compact convex set $\mathcal K$, strict separation would provide a linear form $\ell$ such that $\ell(\P)<\min_{\Q\in\mathcal K}\ell(\Q)$. Proposition~\ref{prop-linear-program-extreme-minimizer} and Theorem~\ref{thm-birkhoff-von-neumann} give
\(\min_{\Q\in\mathcal B_n}\ell(\Q) =\min_{\sigma\in\Perm(n)}\ell(\P_\sigma) =\min_{\Q\in\mathcal K}\ell(\Q),\)
contradicting $\P\in\mathcal B_n$. Hence $\mathcal B_n=\mathcal K$.

The $2n$ row- and column-sum equations have rank $2n-1$, and the strictly positive matrix $\ones_n\transp{\ones_n}/n$ lies in their affine space. Therefore $\mathcal B_n$ has affine dimension $n^2-(2n-1)=(n-1)^2$. Carath\'eodory's theorem in this affine hull now reduces any convex decomposition to at most $(n-1)^2+1$ permutation matrices; zero coefficients can be discarded.
\end{proof}
\(s|I| =\sum_{i\in I}\sum_{j\in N(I)}\R_{i,j} \leq\sum_{j\in N(I)}\sum_i\R_{i,j} =s|N(I)|.\)
Thus $|N(I)|\geq |I|$, so $G_{\R}$ contains a perfect matching.

\(\widehat\P' \eqdef \frac{\R-\lambda\P_\sigma}{s-\lambda}\)
If the current support graph has $E$ edges, the elementary matching subroutine performs at most $n$ breadth-first searches and costs $O(nE)$. Thus a dense implementation costs $O(n^3)$ per extracted permutation and, using the term bound above, at most $O(n^5)$ arithmetic and graph operations overall, with $O(n^2)$ memory.

\paragraph{Rational weights.}
\begin{proposition}[Rational weights as duplicated uniform matching]\label{prop-rational-weights-duplicated-matching}
Let
\(\al=\sum_{i=1}^n \frac{k_i}{N}\delta_{x_i}, \qquad \be=\sum_{j=1}^m \frac{\ell_j}{N}\delta_{y_j}, \qquad \sum_i k_i=\sum_j\ell_j=N,\)
with positive integers $k_i,\ell_j$. Replace each $x_i$ by $k_i$ identical copies and each $y_j$ by $\ell_j$ identical copies, producing two uniform $N$-point clouds. Then the duplicated assignment problem and the discrete Kantorovich problem between $\al$ and $\be$ have the same optimal value. Moreover, there exists an optimal coupling of the form
\(\P_{i,j}=\frac{n_{i,j}}{N}, \qquad n_{i,j}\in\NN, \qquad \sum_j n_{i,j}=k_i, \quad \sum_i n_{i,j}=\ell_j,\)
and couplings whose scaled entries $N\P_{i,j}$ are integral are exactly the collapsed assignments between duplicated clouds. Fractional optimal couplings may nevertheless coexist when the optimum is degenerate.
\end{proposition}
\begin{proof}
Any assignment between the duplicated source and target clouds defines integers $n_{i,j}$ counting how many copied particles of type $x_i$ are matched to copied particles of type $y_j$. These counts satisfy $\sum_j n_{i,j}=k_i$ and $\sum_i n_{i,j}=\ell_j$, and the associated coupling $\P_{i,j}=n_{i,j}/N$ has marginals $k_i/N$ and $\ell_j/N$. The assignment cost is $\frac1N\sum_{i,j} n_{i,j}c(x_i,y_j)=\sum_{i,j}\P_{i,j}c(x_i,y_j)$.
Conversely, any nonnegative integer count matrix with those row and column sums can be realized by allocating the $k_i$ copies of each $x_i$ among the target copies according to $(n_{i,j})_j$. Scale a feasible coupling by $N$ and write $\Q=N\P$. The constraints on $\Q$ have integer right-hand sides $(k_1,\ldots,k_n,\ell_1,\ldots,\ell_m)$. After multiplying the target rows by $-1$, their coefficient matrix is the oriented node-edge incidence matrix of a bipartite graph and is therefore totally unimodular. Every vertex of this transportation polytope is integral. Since a linear objective attains its minimum at a vertex, an integral optimal $\Q$ exists, proving equality of the two optimal values. The argument establishes existence, not integrality of every optimizer; convex combinations of distinct integral optima give fractional optima.
\end{proof}
\paragraph{Applications of discrete transport.}
\subsection{Linear-Programming Algorithms}
\label{sec-kantorovich-lp-algorithms}
\paragraph{Transportation simplex and network simplex.}
\paragraph{Interior-point methods.}
\eql{\label{eq-transport-log-barrier}
\P_\epsilon
\eqdef
\uargmin{\substack{\P\ones_m=\a,\ \P^\top\ones_n=\b\\ \P_{i,j}>0}}
\dotp{\C}{\P}
-
\epsilon\sum_{i,j}\log \P_{i,j},
}
where $\epsilon>0$ is decreased along the algorithm. The barrier is singular at the boundary, so each iterate stays strictly inside the transportation polytope; as $\epsilon\downarrow0$, the central path approaches the set of LP minimizers.

\subsection{Relaxation for Arbitrary Measures}
\label{sec-kantorovich-continuous}
\paragraph{Continuous couplings.}
\begin{defn}[Marginals of a joint measure]\label{def-joint-marginals}
Let $\pi\in\Mm_+^1(\Xx\times\Yy)$ and let $P_\Xx(x,y)=x$, $P_\Yy(x,y)=y$ be the coordinate projections. The marginals of $\pi$ are
\(\pi_1\eqdef(P_\Xx)_\sharp\pi\in\Mm_+^1(\Xx), \qquad \pi_2\eqdef(P_\Yy)_\sharp\pi\in\Mm_+^1(\Yy).\)
Equivalently, for all bounded continuous test functions $f$ on $\Xx$ and $g$ on $\Yy$, \(\int_{\Xx\times\Yy} f(x)\d\pi(x,y)=\int_\Xx f\d\pi_1, \qquad \int_{\Xx\times\Yy} g(y)\d\pi(x,y)=\int_\Yy g\d\pi_2.\)
\end{defn}
\(\int_\Yy \d\pi(x,y)=\d\al(x), \qquad \int_\Xx \d\pi(x,y)=\d\be(y),\)
which is made rigorous by Definition~\ref{def-joint-marginals}, or equivalently by the identities $\pi(A\times\Yy)=\al(A)$ and $\pi(\Xx\times B)=\be(B)$ for measurable sets $A\subset\Xx$ and $B\subset\Yy$.

\begin{defn}[Couplings]\label{def-continuous-couplings}
Given $\al\in\Mm_+^1(\Xx)$ and $\be\in\Mm_+^1(\Yy)$, the set of couplings between $\al$ and $\be$ is
\eql{\label{eq-coupling-generic}
\Couplings(\al,\be)
\eqdef
\enscond{\pi\in\Mm_+^1(\Xx\times\Yy)}{\pi_1=\al \qandq \pi_2=\be}.
}
This is the continuous analogue of the transportation polytope~\eqref{eq-discr-couplings}.
\end{defn}
\begin{defn}[Tensor product and trivial coupling]\label{def-tensor-product-coupling}
Given $\al\in\Mm_+^1(\Xx)$ and $\be\in\Mm_+^1(\Yy)$, the tensor product coupling $\al\otimes\be$ is the probability measure on $\Xx\times\Yy$ defined by
\[
\int_{\Xx\times\Yy} h(x,y)\d(\al\otimes\be)(x,y)
=
\int_\Xx\left(\int_\Yy h(x,y)\d\be(y)\right)\d\al(x)
\]
for every bounded measurable $h$. It is also called the trivial coupling because it makes the two coordinates independent.
\end{defn}
\[
\int_{\Xx\times\Yy} f(x)\d(\al\otimes\be)(x,y)
=
\left(\int_\Xx f(x)\d\al(x)\right)\left(\int_\Yy\d\be(y)\right)
=
\int f\d\al,
\]
\begin{prop}[Product optimality is degenerate]\label{prop-product-coupling-degenerate}
Assume that $\Xx$ and $\Yy$ are compact metric spaces and that $c\in\Cc(\Xx\times\Yy)$.
\(\al\otimes\be\in\uargmin{\pi\in\Couplings(\al,\be)}\int c\,\d\pi, \qquad \text{every coupling in }\Couplings(\al,\be)\text{ is optimal}.\)
They are also equivalent to the additive decomposition of the cost on the product support, \(c(x,y)=u(x)+v(y), \qquad u\in\Cc(\supp(\al)),\quad v\in\Cc(\supp(\be)).\)
\end{prop}
\begin{proof}
If every coupling is optimal, then $\al\otimes\be$ is optimal. Conversely, assume that $\al\otimes\be$ is optimal. We first show that, for every $x_0,x_1\in\supp(\al)$ and $y_0,y_1\in\supp(\be)$,
\(c(x_0,y_0)+c(x_1,y_1) = c(x_0,y_1)+c(x_1,y_0).\)
Indeed, if this equality failed, after exchanging $y_0$ and $y_1$ if necessary one would have a strict inequality
\(c(x_0,y_0)+c(x_1,y_1) > c(x_0,y_1)+c(x_1,y_0).\)
Failure of the equality forces $x_0\neq x_1$ and $y_0\neq y_1$. By continuity, the strict inequality persists with a uniform margin on small neighborhoods $U_0,U_1$ of $x_0,x_1$ and $V_0,V_1$ of $y_0,y_1$, chosen disjoint within each pair. Since the four points lie in the supports, these neighborhoods have positive marginal mass. Denote by $\al_i$ and $\be_i$ the normalized restrictions of $\al$ to $U_i$ and of $\be$ to $V_i$, and choose
\(0<\lambda\leq \min\{\al(U_0)\be(V_0),\al(U_1)\be(V_1)\}.\)
The exchanged measure
\(\tilde\pi = \al\otimes\be -\lambda\,\al_0\otimes\be_0 -\lambda\,\al_1\otimes\be_1 +\lambda\,\al_0\otimes\be_1 +\lambda\,\al_1\otimes\be_0\)
is nonnegative and has the same two marginals as $\al\otimes\be$. The uniform strict inequality on the neighborhoods implies that $\int c\,\d\tilde\pi<\int c\,\d(\al\otimes\be)$, contradicting optimality.

Fixing any $x_\star\in\supp(\al)$ and $y_\star\in\supp(\be)$, the equality of cross differences gives, for all $(x,y)\in\supp(\al)\times\supp(\be)$, \(c(x,y)=c(x,y_\star)+c(x_\star,y)-c(x_\star,y_\star).\)
Thus $c=u+v$ on the product support, with $u(x)=c(x,y_\star)-c(x_\star,y_\star)$ and $v(y)=c(x_\star,y)$. These functions are continuous on the corresponding supports. Every coupling is concentrated on this product support, so for any $\pi\in\Couplings(\al,\be)$,
\(\int c\,\d\pi = \int u\,\d\al+\int v\,\d\be\),
which depends only on the marginals. Hence all couplings are optimal.
\end{proof}
The tensor product is therefore a trivial feasible coupling, not a typical optimizer. Product optimality means that the cost cannot distinguish between dependences once the marginals are fixed.

If there exists a map $T:\Xx\to\Yy$ such that $T_\sharp\al=\be$, then the Monge map induces the graph coupling $\pi=(\Id,T)_\sharp\al\in\Couplings(\al,\be)$, characterized by

\(\int_{\Xx\times\Yy} h(x,y)\d\pi(x,y) = \int_\Xx h(x,T(x))\d\al(x).\)
Applying this identity to $h(x,y)=f(x)$ or $h(x,y)=g(y)$ gives respectively $\pi_1=\al$ and $\pi_2=\be$. Thus graph couplings are precisely the Kantorovich representation of deterministic Monge maps.

\paragraph{Continuous Kantorovich problem.}
\begin{defn}[Continuous Kantorovich problem]\label{def-continuous-kantorovich-problem}
For $\al\in\Mm_+^1(\Xx)$, $\be\in\Mm_+^1(\Yy)$, and a nonnegative Borel cost $c:\Xx\times\Yy\to[0,+\infty]$, define
\eql{\label{eq-mk-generic}
\MK_c(\al,\be)
\eqdef
\inf_{\pi\in\Couplings(\al,\be)}
\int_{\Xx\times\Yy} c(x,y)\d\pi(x,y).
}
A minimizer is an optimal coupling, or optimal transport plan.
The value is understood in $[0,+\infty]$.
\end{defn}
\begin{prop}[Convexity in the marginals and concavity in the cost]\label{prop-kantorovich-value-curvature}
Fix a nonnegative Borel cost $c$. The extended-valued map
\((\al,\be)\longmapsto\MK_c(\al,\be)\)
is jointly convex on pairs of probability measures. Conversely, for fixed marginals $(\al,\be)$, the map $c\mapsto\MK_c(\al,\be)$ is concave on the cone of nonnegative Borel costs: for any such $c_0,c_1$ and $t\in[0,1]$,
\(\MK_{(1-t)c_0+tc_1}(\al,\be) \geq (1-t)\MK_{c_0}(\al,\be)+t\MK_{c_1}(\al,\be).\)
The inequality is understood in $[0,+\infty]$; in particular, the map is concave on every convex family of costs on which it is finite.
\end{prop}
\begin{proof}
For joint convexity, let $(\al_0,\be_0)$ and $(\al_1,\be_1)$ be two pairs of probability measures. The claim is immediate at $t\in\{0,1\}$ or if one of the two values on the right-hand side is infinite. Otherwise, for $\eta>0$, choose $\pi_i\in\Couplings(\al_i,\be_i)$ such that $\int c\d\pi_i\leq\MK_c(\al_i,\be_i)+\eta$ for $i\in\{0,1\}$.
Then $(1-t)\pi_0+t\pi_1$ couples $(1-t)\al_0+t\al_1$ and $(1-t)\be_0+t\be_1$, and hence
\(\MK_c\big((1-t)\al_0+t\al_1,(1-t)\be_0+t\be_1\big) \leq (1-t)\MK_c(\al_0,\be_0)+t\MK_c(\al_1,\be_1)+\eta.\)
Letting $\eta\to0$ proves joint convexity.

For concavity, the endpoint cases are immediate, so assume $t\in(0,1)$ and fix $\pi\in\Couplings(\al,\be)$. Linearity of integration for nonnegative functions gives
\(\int\big((1-t)c_0+tc_1\big)\d\pi \geq (1-t)\MK_{c_0}(\al,\be)+t\MK_{c_1}(\al,\be).\)
Taking the infimum over $\pi$ proves the stated inequality.
\end{proof}
\begin{prop}[Existence for lower-semicontinuous costs]\label{prop-kantorovich-existence-compact}
Assume that $\Xx$ and $\Yy$ are Polish spaces and that $c:\Xx\times\Yy\to[0,+\infty]$ is lower semicontinuous. Then the Kantorovich problem~\eqref{eq-mk-generic} admits at least one minimizer. In particular, this applies to every continuous nonnegative cost on compact metric spaces.
\end{prop}
\begin{proof}
The constraint set is non-empty because it contains the product coupling $\al\otimes\be$. It is uniformly tight: given $\varepsilon>0$, choose compact sets $K_\Xx\subset\Xx$ and $K_\Yy\subset\Yy$ such that $\al(K_\Xx)\geq1-\varepsilon/2$ and $\be(K_\Yy)\geq1-\varepsilon/2$. Every $\pi\in\Couplings(\al,\be)$ then satisfies
\(\pi(K_\Xx\times K_\Yy)\geq1-\varepsilon\).
By Prokhorov's theorem, the set of couplings is relatively compact for weak convergence; the weak topology is reviewed in Section~\ref{sec-wasserstein-topology-applications}. It is also weakly closed, because the coordinate projections are continuous and therefore preserve the marginal identities in the limit. Hence $\Couplings(\al,\be)$ is weakly compact. Finally, Portmanteau's theorem makes $\pi\mapsto\int c\d\pi$ weakly lower semicontinuous. The direct method gives a minimizer.
\end{proof}
\paragraph{Axiomatic characterization of the relaxation.}
\begin{thm}[Kantorovich cost as the maximal convex relaxation]\label{thm-kantorovich-maximal-convex-relaxation}
Assume that $\Xx$ and $\Yy$ are compact metric spaces and that $c:\Xx\times\Yy\to[0,+\infty)$ is continuous. Then $\MK_c$ is the pointwise largest proper functional
\(\mathsf F:\Mm_+^1(\Xx)\times\Mm_+^1(\Yy)\longrightarrow(-\infty,+\infty]\)
that is jointly convex, weakly lower semicontinuous, and satisfies \(\mathsf F(\de_x,\de_y)\leq c(x,y) \qquad\text{for every }(x,y)\in\Xx\times\Yy.\)
More precisely, every such $\mathsf F$ satisfies $\mathsf F(\al,\be)\leq\MK_c(\al,\be)$ for all $(\al,\be)$,
while $\MK_c$ itself has these properties and obeys $\MK_c(\de_x,\de_y)=c(x,y)$.
\end{thm}
\begin{proof}
Joint convexity follows from Proposition~\ref{prop-kantorovich-value-curvature}. For weak lower semicontinuity, consider converging marginals and optimal plans along a subsequence realizing the lower limit. Compactness gives a further weakly convergent subsequence; its limit has the limiting marginals, and lower semicontinuity of the integral gives the claim, exactly as in Proposition~\ref{prop-kantorovich-existence-compact}. The only coupling of $(\de_x,\de_y)$ is $\de_{(x,y)}$, which proves the identity on Dirac pairs.

Fix $\pi\in\Couplings(\al,\be)$. Its marginal constraints are exactly the barycentric identity $(\al,\be)=\int(\de_x,\de_y)\d\pi$.
Jensen's inequality for lower-semicontinuous convex functionals therefore gives \(\mathsf F(\al,\be) \leq\int\mathsf F(\de_x,\de_y)\d\pi(x,y) \leq\int c(x,y)\d\pi(x,y).\)
Taking the infimum over $\pi\in\Couplings(\al,\be)$ proves maximality.
\end{proof}
\paragraph{Monge--Kantorovich equivalence.}
\begin{cor}[Monge--Kantorovich equivalence under Brenier]\label{cor-monge-kantorovich-brenier}
Let $\al,\be\in\Pp_2(\RR^d)$, assume that $\al$ is absolutely continuous with respect to Lebesgue measure, and let $c(x,y)=\norm{x-y}^2$. If $T$ is the Brenier map solving Monge's problem, then $\pi=(\Id,T)_\sharp\al$ is the unique optimal coupling solving the Kantorovich problem. In particular, Monge and Kantorovich costs are the same.
\end{cor}
\begin{proof}
The proof of Brenier's theorem shows that the support of any optimal Kantorovich plan is contained in the subdifferential $\partial\phi$ of a convex function $\phi$. When $\al$ has a density, $\phi$ is differentiable $\al$-almost everywhere, so $\partial\phi(x)=\{\nabla\phi(x)\}$ for $\al$-almost every $x$. Thus every optimal coupling is concentrated on the graph of $T=\nabla\phi$ and must equal $(\Id,T)_\sharp\al$. The graph coupling is feasible and optimal, and the two formulations have the same value.
\end{proof}
\paragraph{Kantorovich solution in 1-D.}\label{sec-1d-kantorovich-solution}
\begin{thm}[One-dimensional Kantorovich solution]\label{prop-1d-kantorovich-quantile-coupling}
Let $\al,\be\in\Mm_+^1(\RR)$ and let $h:\RR\to[0,+\infty)$ be convex. Consider the cost $c(x,y)=h(x-y)$. Write $q_\al\eqdef\cumul{\al}^{-1}$ and $q_\be\eqdef\cumul{\be}^{-1}$, and assume that $h(q_\al-q_\be)\in L^1(0,1)$. Then the quantile coupling
\(\pi^\star=(q_\al,q_\be)_\sharp\mathrm{Leb}_{[0,1]} = (\cumul{\al}^{-1},\cumul{\be}^{-1})_\sharp\mathrm{Leb}_{[0,1]}\)
solves the Kantorovich problem~\eqref{eq-mk-generic} for $c(x,y)=h(x-y)$.
In particular, $h(t)=|t|^p$ gives the usual one-dimensional optimal coupling for every $p\geq1$.
\end{thm}
\begin{proof}
The push-forward statement $\pi^\star\in\Couplings(\al,\be)$ follows from Proposition~\ref{prop-quantile-pushforward}.

The key point is the one-dimensional uncrossing inequality. If $x<x'$ and $y>y'$, set $a=x-y$, $\delta=x'-x>0$ and $\eta=y-y'>0$. Convexity of $h$ implies that increments are monotone, hence
\(h(a+\eta)-h(a) \leq h(a+\delta+\eta)-h(a+\delta),\)
which is exactly
\(h(x-y)+h(x'-y') \geq h(x-y')+h(x'-y)\).
Thus removing a crossing never increases the cost. For a finite transport matrix on two ordered grids, if $i<i'$ and $j>j'$ carry crossed masses $\P_{i,j}$ and $\P_{i',j'}$, one may move $\theta=\min(\P_{i,j},\P_{i',j'})$ units of mass from the crossed entries $(i,j)$, $(i',j')$ to the uncrossed entries $(i,j')$, $(i',j)$. The row and column sums are unchanged, and the preceding inequality shows that the cost does not increase. Repeating this elementary uncrossing step yields an ordered plan; on an ordered uniform quantile grid, this is the diagonal plan.

For general measures, apply the same argument in quantile coordinates. Let $\pi\in\Couplings(\al,\be)$. Let $U_\al$ and $U_\be$ be uniform random variables on $(0,1)$, and denote by $K_\al(x,\d r)$ and $K_\be(y,\d s)$ regular conditional laws of $U_\al$ given $q_\al(U_\al)=x$ and of $U_\be$ given $q_\be(U_\be)=y$. Given $(X,Y)\sim\pi$, draw conditionally $R\sim K_\al(X,\cdot)$ and $S\sim K_\be(Y,\cdot)$. The law $\gamma$ of $(R,S)$ has uniform marginals and satisfies $\pi=(q_\al,q_\be)_\sharp\gamma$.

Let $\kappa_M(r)=\max(-M,\min(r,M))$ and set $q_{\al,M}=\kappa_M\circ q_\al$ and $q_{\be,M}=\kappa_M\circ q_\be$. These functions are bounded and nondecreasing. Approximate them almost everywhere by nondecreasing step functions $q_{\al,M}^N$ and $q_{\be,M}^N$ that are constant on the uniform intervals $I_k=((k-1)/N,k/N]$. The matrix
\(G_{k,\ell}^N=\gamma(I_k\times I_\ell)\)
is a coupling of the two uniform histograms. Proposition~\ref{prop-1d-weighted-sweep}, applied to the ordered step values, therefore gives
\(\int_0^1 h(q_{\al,M}^N(r)-q_{\be,M}^N(r))\d r \leq \int_{(0,1)^2}h(q_{\al,M}^N(r)-q_{\be,M}^N(s))\d\gamma(r,s).\)
For fixed $M$, the arguments of $h$ remain in the compact interval $[-2M,2M]$. Dominated convergence as $N\to\infty$ yields the same inequality for $q_{\al,M}$ and $q_{\be,M}$.

Finally, $\kappa_M$ is nondecreasing and $1$-Lipschitz. Hence, for every $x,y\in\RR$, there exists $t_M(x,y)\in[0,1]$ such that
\(\kappa_M(x)-\kappa_M(y)=t_M(x,y)(x-y)\).
Convexity and nonnegativity imply
\(h(t_M(x,y)(x-y)) \leq (1-t_M(x,y))h(0)+t_M(x,y)h(x-y) \leq h(0)+h(x-y).\)
The assumed integrability controls the diagonal term. If the cost of $\pi$ is finite, the same bound controls the right-hand side; if it is infinite, the desired comparison is immediate. Dominated convergence as $M\to\infty$ therefore gives
\(\int_0^1 h(q_\al(r)-q_\be(r))\d r \leq \int h(x-y)\d\pi(x,y)\)
for every $\pi\in\Couplings(\al,\be)$.
\end{proof}
\subsection{Cone of Convex Functions and Monge Gap}
\label{sec-monge-gap}
\paragraph{Definition and optimality certificate.}
\begin{defn}[Monge gap]\label{def-monge-gap}
Let $\al\in\Mm_+^1(\Xx)$, let $T:\Xx\to\Yy$ be measurable, and assume that the following costs are finite. The $c$-Monge gap of $T$ relative to $\al$ is
\eql{\label{eq-monge-gap}
\mathcal M_\al^c(T)
\eqdef
\int c(x,T(x))\d\al(x)-\MK_c(\al,T_\sharp\al).
}
\end{defn}
\begin{prop}[Elementary properties of the Monge gap]\label{prop-monge-gap-properties}
Under the standing Polish-space assumptions, suppose that the relevant integrals are finite. Then $\mathcal M_\al^c(T)\geq0$, with equality if and only if $(\Id,T)_\sharp\al$ is an optimal coupling between $\al$ and $T_\sharp\al$. Moreover,
\eql{\label{eq-monge-gap-self-coupling}
\mathcal M_\al^c(T)
=
\sup_{\eta\in\Couplings(\al,\al)}
\int\big[c(x,T(x))-c(x,T(z))\big]\d\eta(x,z).
}
Finally, the gap is unchanged when $c(x,y)$ is replaced by $c(x,y)+r(x)+s(y)$, provided the modified cost remains admissible and all marginal terms are integrable.
\end{prop}
\begin{proof}
The graph plan $(\Id,T)_\sharp\al$ is feasible for $\MK_c(\al,T_\sharp\al)$, which proves nonnegativity and the equality criterion. Every $\eta\in\Couplings(\al,\al)$ induces the coupling $(x,z)\mapsto(x,T(z))$. Conversely, disintegrating $\al$ along $T$ lifts every coupling in $\Couplings(\al,T_\sharp\al)$ to such a self-coupling. Subtracting the resulting infimum from the graph cost gives~\eqref{eq-monge-gap-self-coupling}. Adding $r(x)+s(y)$ shifts both terms in~\eqref{eq-monge-gap} by $\int r\d\al+\int s\d(T_\sharp\al)$.
\end{proof}
\paragraph{Convex gaps and convex potentials.}
\begin{prop}[Convex Monge gaps for affine-difference costs]\label{prop-affine-difference-monge-gap}
Let $\Yy\subset\RR^m$ be convex and let the nonnegative cost $c$ have the form
\eql{\label{eq-affine-difference-cost}
c(x,y)=a(x)+b(y)-\dotp{A(x)}{y},
}
where $A,T\in L^2(\al;\RR^m)$ and all displayed terms are finite. Then
\eql{\label{eq-affine-difference-monge-gap}
\mathcal M_\al^c(T)
=
\sup_{\eta\in\Couplings(\al,\al)}
\int\dotp{A(x)}{T(z)-T(x)}\d\eta(x,z).
}
Moreover, $\mathcal M_\al^c(T)=0$ if and only if there is a proper lower-semicontinuous convex function $\varphi$ such that
\eql{\label{eq-affine-difference-subgradient-characterization}
T(x)\in\partial\varphi(A(x))
\qquad\text{for $\al$-almost every }x.
}
\end{prop}
\begin{proof}
The marginal terms $a$ and $b$ cancel in~\eqref{eq-monge-gap-self-coupling}, giving~\eqref{eq-affine-difference-monge-gap}; this is a supremum of linear functionals of $T$. Next, disintegrating $\al$ along $A$ lifts every coupling between $A_\sharp\al$ and $T_\sharp\al$. Hence zero gap is equivalent to optimality of $(A,T)_\sharp\al$ for the bilinear cost $(u,y)\mapsto-\dotp{u}{y}$. Up to marginal terms, this is the quadratic cost $\frac12\norm{u-y}^2$. By Rockafellar's characterization, recalled in Section~\ref{sec-cyclical-monotonicity}, its optimal plans are exactly those concentrated on the subdifferential of a proper lower-semicontinuous convex function.
\end{proof}
\(B_\Phi(y|x)=\Phi(y)-\Phi(x)-\dotp{\nabla\Phi(x)}{y-x},\)
where $\Phi$ is differentiable and convex: here $a(x)=\dotp{\nabla\Phi(x)}{x}-\Phi(x)$, $b(y)=\Phi(y)$ and $A(x)=\nabla\Phi(x)$. The argument order is essential: on a full-dimensional domain, the reversed divergence $B_\Phi(x|y)$ is generally not of the form~\eqref{eq-affine-difference-cost} unless $\Phi$ is quadratic.

\paragraph{One-dimensional order and isotonicity.}
\begin{defn}[Increasing rearrangement relative to a source]\label{def-increasing-rearrangement-map}
Let $\al\in\Mm_+^1(\RR)$ be atomless and let $T:\RR\to\RR$ be measurable. The increasing rearrangement of $T$ relative to $\al$ is
\eql{\label{eq-increasing-rearrangement-map}
T^{\uparrow,\al}
\eqdef
\cumul{T_\sharp\al}^{-1}\circ\cumul{\al}.
}
It is the unique $\al$-almost-everywhere nondecreasing map satisfying $(T^{\uparrow,\al})_\sharp\al=T_\sharp\al$.
\end{defn}
\begin{prop}[One-dimensional Monge-gap formulas]\label{prop-monge-gap-one-dimensional}
Let $\al\in\Mm_+^1(\RR)$ be atomless, let $c(x,y)=h(x-y)$ with $h$ convex, and assume the relevant terms are integrable. Then
\eql{\label{eq-monge-gap-increasing-rearrangement}
\mathcal M_\al^c(T)
=
\int\big[h(x-T(x))-h(x-T^{\uparrow,\al}(x))\big]\d\al(x).
}
If $h$ is strictly convex, this gap vanishes if and only if $T$ has a nondecreasing representative $\al$-almost everywhere.
For the equal-weight atomic measure $\al_n=\frac1n\sum_{i=1}^n\de_{x_i}$ on the uniform grid $x_i=x_1+(i-1)\Delta$, let $t_i=T(x_i)$ and let $t_{(1)}\leq\cdots\leq t_{(n)}$ be their order statistics. For $c(x,y)=|x-y|^2$,
\eql{\label{eq-monge-gap-uniform-grid}
\mathcal M_{\al_n}^{|\cdot|^2}(T)
=
\frac1n\sum_{i=1}^n\big[(x_i-t_i)^2-(x_i-t_{(i)})^2\big]
=
\frac{2\Delta}{n}\sum_{1\leq i<j\leq n}(t_i-t_j)_+.
}
\end{prop}
\begin{proof}
The one-dimensional rearrangement result of Theorem~\ref{prop-1d-quantile-map}, together with its coupling extension in Theorem~\ref{prop-1d-kantorovich-quantile-coupling}, shows that $T^{\uparrow,\al}$ is optimal between $\al$ and $T_\sharp\al$. If $h$ is strictly convex, atomlessness makes this optimal map unique $\al$-almost everywhere, so zero gap is equivalent to $T=T^{\uparrow,\al}$ $\al$-almost everywhere. For the empirical formula, sorting $(t_i)_i$ gives the increasing optimal assignment and hence the first equality. An adjacent exchange of an inversion $u>v$ decreases the average squared cost by $2\Delta(u-v)/n$; bubble sort exchanges every inverted pair once, which gives the second equality.
\end{proof}
\paragraph{Monge-gap-regularized kernel regression.}\label{par-monge-gap-kernel-regression}
The hinge representation above turns transport optimality into a tractable soft monotonicity prior. Let $x_i=x_1+(i-1)\Delta$ and let $y_i\in\RR$ be noisy observations.

\eql{\label{eq-monge-gap-kernel-ridge}
\umin{T\in\RKHS_k}
\left\{
\frac1n\sum_{i=1}^n|T(x_i)-y_i|^2
+\lambda\norm{T}_{\RKHS_k}^2
+\gamma\mathcal M_{\al_n}^{|\cdot|^2}(T)
\right\},
\qquad \lambda>0,\quad \gamma\geq0,
}
where $\al_n=n^{-1}\sum_i\de_{x_i}$. The usual orthogonal-projection argument shows that a minimizer has the representer form $T_\q(\cdot)=\sum_{j=1}^n\q_j k(x_j,\cdot)$.

\eql{\label{eq-monge-gap-kernel-finite}
\umin{\q\in\RR^n}
\left\{
\frac1n\norm{\K\q-\mathrm y}^2
+\lambda\dotp{\q}{\K\q}
+\frac{2\gamma\Delta}{n}
\sum_{1\leq i<j\leq n}
\big[(\K\q)_i-(\K\q)_j\big]_+
\right\}.
}
This is a finite-dimensional convex program: the first two terms are convex quadratics because $\K\succeq0$, and the last term is a sum of hinges of affine functionals. Thus every point along the regularization path can be computed globally, without alternating between regression and transport.

\subsection{\texorpdfstring{$c$}{c}-Cyclical Monotonicity}
\label{sec-cyclical-monotonicity}
\paragraph{\texorpdfstring{Support and $c$-cyclical monotonicity.}{Support and c-cyclical monotonicity.}}
The support of a coupling is the topological support introduced in Definition~\ref{def:support}, now applied to a Radon measure on $\Xx\times\Yy$. Thus $(x,y)\in\supp(\pi)$ exactly when every open neighborhood of $(x,y)$ has positive $\pi$-mass.

\begin{defn}[$c$-cyclical monotonicity]\label{def:ccm}
A set $\Gamma\subset\Xx\times\Yy$ is $c$-cyclically monotone if, for every $k\geq2$, every finite family $(x_i,y_i)_{i=1}^k\subset\Gamma$ and every permutation $\sigma$ of $\{1,\ldots,k\}$,
\(\sum_{i=1}^k c(x_i,y_i) \leq \sum_{i=1}^k c(x_i,y_{\sigma(i)}).\)
\end{defn}
\(\sum_{i=1}^k c(x_i,y_i) \leq \sum_{i=1}^k c(x_i,y_{i+1}), \qquad y_{k+1}=y_1.\)
\paragraph{Optimal matching to optimal transport.}
Let the marginals be uniform on $n$ points, $\al=\frac1n\sum_{i=1}^n\delta_{x_i}$ and $\be=\frac1n\sum_{i=1}^n\delta_{y_i}$. By Corollary~\ref{cor-kantorovich-matching}, there exists an optimal plan induced by a permutation. Its support $\Gamma=\{(x_i,y_{\sigma(i)})\}_i$ is $c$-cyclically monotone: otherwise exchanging the finitely many targets along a violating cycle would lower the matching cost.

\begin{thm}[Optimal plans are $c$-cyclically monotone]\label{thm:opt_ccm}
Assume that $c:\Xx\times\Yy\to[0,+\infty)$ is continuous and that the Kantorovich value is finite. For any optimal plan $\pi$ solving~\eqref{eq-mk-generic}, $\supp(\pi)$ is $c$-cyclically monotone.
\end{thm}
\begin{proof}
Suppose that $\supp(\pi)$ is not $c$-cyclically monotone. Then there exist points $(x_i,y_i)_{i=1}^k$ in the support and a permutation $\sigma$ such that $\sum_i c(x_i,y_i)>\sum_i c(x_i,y_{\sigma(i)})$.
By continuity of $c$, after shrinking neighborhoods $U_i\ni x_i$ and $V_i\ni y_i$, the same strict inequality holds uniformly for every choice of points in these neighborhoods:
\(\sum_i c(u_i,v_i) > \sum_i c(u_i,\tilde v_{\sigma(i)}) \qquad (u_i\in U_i,\ v_i\in V_i,\ \tilde v_{\sigma(i)}\in V_{\sigma(i)}).\)
Choose the sets so that $m_i\eqdef\pi(U_i\times V_i)>0$. Taking
\(0<\lambda\leq\left(\sum_{i=1}^k\frac1{m_i}\right)^{-1}\)
ensures that the scaled restrictions
\(\pi_i=\lambda\frac{\pi|_{U_i\times V_i}}{m_i}\)
have common mass $\lambda$ and satisfy $\sum_i\pi_i\leq\pi$, even if some rectangles overlap. Let $\al_i=(P_\Xx)_\sharp\pi_i$ and $\be_i=(P_\Yy)_\sharp\pi_i$, and define $\tilde\pi=\pi-\sum_i\pi_i+\sum_i(\al_i\otimes\be_{\sigma(i)})/\lambda$.
The removed and reinserted first marginals are both $\sum_i\al_i$, and the removed and reinserted second marginals are both $\sum_i\be_i$ because $\sigma$ is a permutation. Hence $\tilde\pi\in\Couplings(\al,\be)$. Integrating the uniform strict inequality against the product probability $\otimes_i(\pi_i/\lambda)$ shows that the reinserted crossed terms have strictly smaller cost than the removed diagonal terms. This contradicts the optimality of $\pi$.
\end{proof}
\paragraph{Monotonicity.}
Assume the optimal plan is induced by a measurable map $T:\Xx\to\Yy$, meaning that $\pi=(\Id,T)_\sharp\al$. There is a set $G\subset\Xx$ of full $\al$-measure such that $(x,T(x))\in\supp(\pi)$ for every $x\in G$. For any $x_1,\ldots,x_k\in G$, cyclical monotonicity reads

\(\sum_{i=1}^k c(x_i,T(x_i)) \leq \sum_{i=1}^k c(x_i,T(x_{i+1})), \qquad x_{k+1}=x_1.\)
\paragraph{One dimension.}
\(|x-T(x)|^p+|y-T(y)|^p \leq |x-T(y)|^p+|y-T(x)|^p,\)
For $p>1$, strict convexity implies that a reversed pair $x<y$ and $T(x)>T(y)$ violates this inequality. Hence every optimal map is nondecreasing on the full-measure transport set, recovering the classical monotone rearrangement.